\chapter{RL Is Not What You Think}
\label{ch:rl-not-what-you-think}

Most introductions to reinforcement learning begin with the Bellman equation, march through Q-learning, and end with a collection of algorithms.
This framing creates a dangerous illusion: that RL is merely supervised learning with delayed labels.
It is not.
The single most important thing to understand about RL is its \emph{sampling structure}---and how that structure differs from supervised learning at a fundamental level.

% ===========================================================
\section{The Same Optimization Framework}
% ===========================================================

At the level of code, RL and supervised learning (SL) are remarkably similar.
Both minimize a loss with gradient descent:
%
\begin{equation}
\theta \leftarrow \theta - \alpha\,\nabla_\theta \mathcal{L}(\theta).
\end{equation}
%
In SL the loss is typically a prediction error:
$\mathcal{L} = L\bigl(f_\theta(x),\, y\bigr)$.
In RL the loss takes one of two canonical forms:
%
\begin{itemize}
    \item \textbf{Value-based:}
        $\mathcal{L} = \bigl(Q_\theta(s,a) - y\bigr)^2$,
        where $y = r + \gamma \max_{a'} Q_{\bar\theta}(s',a')$.
    \item \textbf{Policy-based:}
        $\mathcal{L} = -\log\pi_\theta(a|s)\,A(s,a)$.
\end{itemize}
%
Both use SGD or Adam.  Both backpropagate through the same autodiff graph.
At the ``write a training loop'' level, RL and SL share the same optimization framework:
forward pass, compute loss, backward pass, update.

\begin{keyidea}[Surface Similarity]
RL and SL share the same \texttt{loss.backward(); optimizer.step()} pipeline.
The difficulty of RL is \emph{not} in the optimizer---it is in where the data comes from.
\end{keyidea}

% ===========================================================
\section{The Closed Loop: Where Everything Changes}
\label{sec:closed-loop}
% ===========================================================

In SL, the dataset $\mathcal{D} = \{(x_i, y_i)\}$ is fixed before training begins.
The model changes; the data does not.
This is an \textbf{open-loop} structure:
%
\begin{equation}
\text{Fixed } \mathcal{D} \;\longrightarrow\; \mathcal{L} \;\longrightarrow\; \nabla_\theta \;\longrightarrow\; \theta' \;\longrightarrow\; \text{better model.}
\end{equation}

In RL, data is generated by the policy itself.
When $\theta$ changes, the policy $\pi_\theta$ changes, which changes the sampling distribution, which changes the data, which changes the loss landscape.
This is a \textbf{closed loop}:
%
\begin{equation}
\pi_\theta \;\longrightarrow\; \text{sample data} \;\longrightarrow\; \mathcal{L} \;\longrightarrow\; \nabla_\theta \;\longrightarrow\; \pi_{\theta'} \;\longrightarrow\; \text{new data distribution.}
\end{equation}

\begin{keyidea}[The Closed Loop]
In SL, the algorithm does not affect the data distribution.
In RL, the algorithm \emph{is} the data distribution generator.
The policy is simultaneously the object being optimized and the source of training data.
\end{keyidea}

An analogy helps sharpen this:
\begin{itemize}
    \item \textbf{SL} is a student with a pre-printed textbook.
          How the student studies does not change the textbook.
    \item \textbf{RL} is a student who writes their own textbook as they learn.
          What they study affects what practice problems appear next,
          which affects what they learn after that.
          A premature bias produces a biased textbook for all future study.
\end{itemize}

% ===========================================================
\section{Four Links in the Causal Chain}
\label{sec:four-links}
% ===========================================================

The closed loop can be decomposed into four causal links, each feeding into the next:
%
\begin{enumerate}
    \item \textbf{Policy determines sampling.}
          The current $\pi_\theta(a|s)$ decides which state--action pairs are visited.
    \item \textbf{Sampling determines data.}
          The visited pairs form the empirical distribution on which all estimates are based.
    \item \textbf{Data determines value estimates.}
          $Q$, $V$, and advantage $A$ are all estimated from the sampled data---they inherit its biases and gaps.
    \item \textbf{Value estimates update the policy.}
          Whether through $\arg\max Q$ or advantage-weighted log-likelihood, the estimated values reshape $\pi_\theta$.
\end{enumerate}
%
Then the cycle repeats.
Every RL-specific technique---replay buffers, target networks, importance sampling, clipping, entropy bonuses---exists to manage pathologies that arise \emph{somewhere in this loop}.

\begin{table}[ht]
\centering
\caption{RL techniques mapped to the closed-loop problems they address.}
\label{tab:closed-loop-fixes}
\begin{tabular}{ll}
\toprule
\textbf{Technique} & \textbf{Closed-Loop Problem Addressed} \\
\midrule
Replay buffer       & Temporal correlation; distribution drift \\
Target network      & Bootstrap target chasing itself \\
$\epsilon$-greedy / entropy bonus & Sampling collapse; insufficient coverage \\
Importance sampling & Distribution mismatch (old vs.\ new policy) \\
PPO clipping        & Excessive policy change per step \\
Double Q-learning   & $\max$ operator amplifying estimation noise \\
\bottomrule
\end{tabular}
\end{table}

None of these techniques are necessary in SL, because SL's data distribution is not affected by the model.

% ===========================================================
\section{RL Labels Are Not Given---They Are Constructed}
% ===========================================================

In SL, labels $y$ are externally provided and independent of the model.
In RL, the ``labels'' are constructed by the learning system itself:
%
\begin{itemize}
    \item In value-based RL, the target $y = r + \gamma \max_{a'} Q_{\bar\theta}(s', a')$ depends on the model's own Q-estimates.
          The model creates its own labels and then fits them---\emph{self-bootstrapping regression}.
    \item In policy gradient methods, the advantage $A(s,a)$ is estimated from sampled returns or a learned critic, both of which depend on the current policy.
\end{itemize}
%
This self-referential structure means that errors compound: a biased Q-estimate produces biased targets, which produce more biased Q-estimates.
Understanding this recursive dependency is essential for understanding why RL is unstable in ways that SL is not.

% ===========================================================
\section{The Real Definition of RL}
\label{sec:real-definition}
% ===========================================================

With the closed loop in mind, we can state what RL truly is---not in textbook terms, but in operational terms:

\begin{keyidea}[RL in One Sentence]
Reinforcement learning is \textbf{biased preference optimization under shifting distributions}:
on a constantly changing data distribution, using biased sampled data,
estimate long-term returns, and optimize the next round's sampling strategy.
\end{keyidea}

This definition highlights three properties that separate RL from SL:
\begin{enumerate}
    \item \textbf{Shifting distributions.}  The data distribution is non-stationary because the policy generates it.
    \item \textbf{Biased estimates.}  Q-values, advantages, and returns are all estimates with sampling noise, bootstrap bias, and function approximation error.
    \item \textbf{Preference optimization.}  The goal is not to find a ``true'' value function in some god's-eye sense, but to stably push the behavioral distribution toward actions that are statistically better.
\end{enumerate}

This is why modern RL has drifted away from the Bellman equation as a primary optimization target and toward methods that directly shape the policy distribution (PPO, SAC, RLHF, GRPO).
The Bellman equation describes an ideal fixed point; the training process is a noisy, nonlinear, recursive iteration that must be stabilized by engineering.

% ===========================================================
\section{Why This Matters for Everything That Follows}
% ===========================================================

The rest of this book traces two consequences of the closed loop:

\begin{enumerate}
    \item \textbf{The macro arc} (Chapter~\ref{ch:grand-arc}): how the field evolved from dynamic programming---where the Bellman equation is everything---to DPO/GRPO, where the Bellman equation is completely absent.
          Each stage is a response to the closed loop's instabilities.
    \item \textbf{The micro mechanisms} (Part~II): how specific techniques (loss design, on/off-policy structure, exploration, bias control, compression, action modeling) each address a particular link in the four-step causal chain.
\end{enumerate}

The unifying theme is simple:
\emph{RL is not hard because the optimization is hard.
RL is hard because the ground shifts under the optimizer's feet.}
Every algorithm in this book is, at its core, a strategy for optimizing on shifting ground.
